\documentclass[12pt]{article}
\usepackage[utf8]{inputenc}
\usepackage[T1]{fontenc}
\usepackage{amsmath, amssymb, geometry, booktabs}
\usepackage{enumitem}
\geometry{margin=1in}
\setlength{\parindent}{0pt}
\renewcommand{\arraystretch}{1.2}

\title{Final Exam Review}
\author{ECON 308 -- International Economic Relations}
\date{Fall 2025}

\begin{document}
\maketitle

% -------------------------------------------------------------

\section*{Problem 1 — Ricardian Model}

Two countries produce wheat (W) and cloth (C). Unit labor requirements:

\begin{center}
\begin{tabular}{lcc}
\toprule
Country & $a_{LW}$ & $a_{LC}$ \\
\midrule
Home & 4 & 2 \\
Foreign & 5 & 1 \\
\bottomrule
\end{tabular}
\end{center}

Home has 200 workers; Foreign has 150.

\begin{enumerate}[label=(\arabic*)]
\item Determine each country's autarky relative price of wheat, $P_W/P_C$.
\item Identify who has comparative advantage in each good.
\item Suppose the world relative price is $P_W/P_C = 1.5$.
    \begin{itemize}
        \item Which good does each country export?
        \item Compute each country's production bundle under free trade.
    \end{itemize}
\item Explain why both countries gain from trade even if one is absolutely less productive.
\end{enumerate}

% -------------------------------------------------------------

\section*{Problem 2 — Specific Factors Model}

Home produces food (F) and manufactures (M). Labor is mobile; capital is specific to M; land is specific to F.

\[
Q_M = 5K^{1/2}L_M^{1/2}, \qquad Q_F = 4T^{1/2}L_F^{1/2}.
\]

Total labor: $L = 200$. Capital: $K = 100$. Land: $T = 100$.

\begin{enumerate}[label=(\arabic*)]
\item Derive the marginal product of labor in each sector.
\item Set up the equilibrium condition for labor allocation.
\item Suppose the relative price of manufactures rises.
    \begin{itemize}
        \item What happens to wages?
        \item What happens to the return to capital and the return to land?
    \end{itemize}
\item Explain who gains and who loses from this price change.
\end{enumerate}

% -------------------------------------------------------------

\section*{Problem 3 — Heckscher--Ohlin Model}

Two goods: machinery (M) is capital-intensive; textiles (T) is labor-intensive.

\begin{center}
\begin{tabular}{lcc}
\toprule
Country & $K$ & $L$ \\
\midrule
Home & 900 & 400 \\
Foreign & 300 & 600 \\
\bottomrule
\end{tabular}
\end{center}

\begin{enumerate}[label=(\arabic*)]
\item Identify which country is capital-abundant and which is labor-abundant.
\item According to Heckscher--Ohlin, who exports which good?
\item Using Stolper--Samuelson, explain who gains and who loses from opening to trade in each country.
\item Using Rybczynski, describe how an increase in Home labor affects output of M and T.
\end{enumerate}

% -------------------------------------------------------------

\section*{Problem 4 — Standard Trade Model}

A country produces apples (A) and electronics (E):

\[
A^2 + 4E^2 = 1600.
\]

World price ratio:
\[
P_E/P_A = 2.
\]

Preferences: $E = 0.25A$.

\begin{enumerate}[label=(\arabic*)]
\item Find the production point that maximizes the value of output.
\item Determine the consumption bundle.
\item Calculate exports and imports.
\item Suppose $P_E/P_A$ rises.
    \begin{itemize}
        \item Explain how production and consumption adjust.
        \item Does welfare rise or fall?
    \end{itemize}
\item Define export-biased growth and explain its effect on the terms of trade.
\end{enumerate}

% -------------------------------------------------------------

\section*{Problem 5 — External Economies of Scale and Location}

Two countries A and B can produce semiconductors. There are external economies of scale, so average cost falls with industry size. A currently hosts the industry; technologies are identical.

\begin{enumerate}[label=(\arabic*)]
\item Explain why A may remain the global semiconductor producer even though B has the same technology.
\item Describe conditions under which B could successfully enter the industry.
\item Suppose B subsidizes semiconductor production.
    \begin{itemize}
        \item Explain how this may shift global production to B.
        \item Discuss welfare implications for A, for B, and globally.
    \end{itemize}
\end{enumerate}

% -------------------------------------------------------------

\section*{Problem 6 — Trade Policy (Tariffs)}

Demand and supply for steel in Home:

\[
Q_D = 800 - 2P, \qquad Q_S = P - 40.
\]

World price: $P_W = 200$.  
Foreign export supply:
\[
P^* = 150 + 0.5Q^*.
\]

\begin{enumerate}[label=(\arabic*)]
\item Find the free-trade price and quantity of imports.
\item Home imposes a tariff of $t = 30$.
    \begin{itemize}
        \item Find the new world price.
        \item Find the new domestic price.
        \item Compute imports.
    \end{itemize}
\item Compute:
    \begin{itemize}
        \item Tariff revenue
        \item Terms-of-trade gain
        \item Deadweight loss
    \end{itemize}
\item Determine whether national welfare rises or falls.
\end{enumerate}

\end{document}
